\documentclass[12pt]{article}

\usepackage[a4paper, total={6in, 8in}]{geometry}
\usepackage{textcomp}

\begin{document}
\setlength{\parindent}{0pt}

\def \t {\textrightarrow}
\def \v {\vspace{1em}}
\def \bi {\begin{itemize}}
\def \ei {\end{itemize}}
\def \s[#1] {\section*{#1}}
\def \ss[#1] {\subsection*{#1}}
\def \sss[#1] {\subsubsection*{#1}}

\s[Pirandello]
La anziana truccata fa avvertire il contrasto \t prima si ride, poi si empatizza \t in questo caso lei viene osservata con occhi amorevole \\
Calvino è l'autore della leggerezza \t che non vuol dire superficialità, e l'iniziatore della leggerezza è Palazzeschi \t elabora il "Manifesto del contro dolore" \\
La sua ilarità nasconde un dolore \t ha dovuto sopprimere la sua omosessualità \\
Lui vive un arco del novecento \t e ha visto una società che ha affermato la consapevolezza che non è il genere a differenziare la persona \\
Leggerezza: chiedere definizione

\v

Dramma è un'opera teatrale, non poema \\
Pirandello è affacciato alla cultura europea, e andrà anche a Bonn \t la sua attenzione al lettore arriva anche dal fatto che pubblica l'intera sua raccolta \\
Mondadori pubblica "Maschere nude" \t spesso accompagnato dai dipindi di De Chirico di persone come manichini \t questo connette la visione dell'artista con la crisi del 900 \\
Novecento \t si parla di interventismo/neutralismo, appoggio o no al fascismo \t manifesto degli intellettuali antifascisti \t Pirandello firma quello degli intellettuali fascisti

\v

Qual'è il problema da cui nascono tutte le sue riflessioni? \t le sue opere sono anche creatività di pensiero, e crea romanzi/saggio (che hanno fitta rete di riflessioni) \t e dialogo continuo tra i personaggi = precursore dell'evoluzione teatrale \\
Il passaggio da romanzo a opera teatrale è immediato \t i dialoghi sono molti \\
Crediamo di essere uno, ma in realtà siam molti \t e quando crediamo di conoscersi, viviamo il dramma di confrontarci con gli altri \\
Ci vediamo da punti di vista diversi \\
Comicità e umorismo sono diversi \t umorismo = andare al di là di una parvenza, imporre una visione democratica di andare oltre, mentre comcità è ridere e fermarsi \\
Gli elementi fondamentali escono con questo

\v

Priandello nasce ad Agrigento nel 67? \t vive la sicilianità evidenziata con Verga, e poi scoppia la zolfatara \\
Ha pero formazione classica \t studia legge e lettere a Palermo, e realizza visione politica che lo fa entrare in conflitto con il padre \\
Questo fa si che l'autore vive una dimensione di sensibilià interiore \\
La moglie ha manie di persecuzione \t che esprimono un animo diviso, in balia di diverse visioni contrastanti \t e partono da eventi drammatici, e non sempre si possono risolvere \\
Nella sua parabola di vita la sua sensibilità lo porta ad avere un occhio allertato verso il dolore \t la moglie pensa di essere tradita da lui, ed è vero \\
La moglie viene assistita dopo il trauma dalla famiglia, che è dissestata \t e la sua stabilità interiore alla fine non è + gestibile \\
La malattia diventa ingestibile dal punto di vista familiare \t e la moglie viene rinchiusa in un manicomio \\
Sensi di colpa \t nella lucidità si conserva \\
La sua produzione converge anche se testi di breve entità, come le novelle \t la sua produzione e amore per la letteratura portano alla sua affermazione universitaria \\
Si laurea a Bonn, e questa scelta è per trovare un luogo per condurre i suoi studi filologici \t qua impara il tedesco, e legge autori che lo influenzano \\
Nella sua parabola universitaria si affianca a molti autori \t come Thomas Mann, Rilke \t con la parabola intimistica \\
La contestazione della politica, è stato firmatario degli intellettuali fascisti \t perchè lui ravvisava un bisogno di ordine nella cultura, e questa immagine di uno stato corrotto appariva ai suoi occhi come rimediabile \\
Non si aspettava però la violenza \t Montale invece firma quello degli intellettuali antifascisti \\
È ancora sposato con la Portulano \\
Formazione radicata nel tempo, dell'attore che si sta costruendo \\
La vita è il teatro \t massima dell'autore \\
Nel 29 diventa senatore? \t nel 34 Nobel per la letteratura \\
Capisce a due anni dalla morte \t lo porta ad allontanarsi dalla decisione che aveva preso di firmare il manifesto \t le sue polveri vengono lasciate nella sua casa \t morte di polmonite

\ss[L'esclusa]
Aiala, Abba, la sicilianita in una dimensione ancora verista \t viene descritto un mondo ancora verghiano \\
L'opera inaugura il romanzo nuvo \t Marta Aiala è una giovane donna (il romanzo viene chiamato cosi, poi viene rinominato) \t l'esclusione, la figura femminile che si spersonalizza \\
Lei non è nel suo mondo, non è se stessa \t gli altri per una sfortuna del destino la vedono per quello che lei potrebbe non essere stata \\
Nella sua famiglia, le figure maschili sono patriarcali \t la donna doveva essere l'"angelo del focolare", accudente, figura materna = non donna a 360 gradi \\
Marta Aiala non diventa madre, ma intrattiene per caso scambi culturali con un uomo che non è il marito \t si prefigura già un triangolo amoroso \\
Parte conclusiva del 900 \t nel 1893 esce marta aiala, che è stato in gestazione fino al 1901 (quindi ha tempo di rivederlo) \\
Con questo romanzo si vede la tratteggiata dell'esclusa, che diventa esclusa nel 1901 \\
La donna si vuole sdoganare culturalmente \t le suffragette, formazione femminile era negata \t ed era negata la parola nei sindacati \\
Questo universo presenta gia innovazione nel romanzo \t lei ha interessi culturali, si vuole formare \t intrattiene uno scambio epistolare con Gregorio Albiniani, mentre è sposata con Rocco Penagori? \\
Il marito trova le lettere, e interpreta le lettere come tradimento \\
Il marito allontana lei senza regolare processo, e con lui tutto il paese \t c'è il personaggio in più che è il paese \\
Marta Aiala si riscatta, studia, vince concorso e quando deve essere immessa in ruolo, viene bandita da tutte le scuole \t finchè non riesce ad entrare in una scuola, dove i genitori sono disaccordo \\
Come in una conchiglia chiusa, la sicilia rimane limitata = limite \\
L'insegnamento le portera progressi \t pero c'p individuo che la importuna, quindi si sposta a Roma \t incontra il mondo politico del tempo, ma il romanzo torna a punto di partenza quando muore la suocera \\
Lei viene chiamata al capezzale, e quando lei rimane incinta viene accolta dal marito che la riaccoglie e perdona
\end{document}
